%%
%% Copyright (c) 2018-2020 Weitian LI <wt@liwt.net>
%% CC BY 4.0 License
%%
%% Created: 2018-04-11
%%

% Chinese version
\documentclass[zh]{resume}

% File information shown at the footer of the last page
\fileinfo{%
  \faCopyright{} 2025--2026, 阿东玩AI \hspace{0.5em}
  \creativecommons{by}{4.0} \hspace{0.5em}
  \githublink{adonogwanai}{LLM-Resume-Template} \hspace{0.5em}
  \faEdit{} \today
}

\name{邓}{祎明}

\keywords{算法工程师, 计算机视觉, 眼动跟踪, 目标检测, 目标跟踪, 端侧部署}

% \tagline{\texorpdfstring{\icon{\faBinoculars} }{}<position-to-look-for>}
% \tagline{<current-position>}

% \photo[
%   shape=<circular|square>,    % default is circular
%   position=<left|right>,      % default is left
% ]{<width>}{<filename>}
% Example:
% \photo[shape=square]{5.5em}{photo}

\profile{
  \mobile{177-7918-0948}
  \email{dengyiming@std.uestc.edu.cn}
  \university{电子科技大学}
  \degree{生物医学工程 \textbullet 硕士}
  \birthday{2002-09}
  % Custom information:
  % \icontext{<icon>}{<text>}
  % \iconlink{<icon>}{<link>}{<text>}
}

\begin{document}
\makeheader

%======================================================================
% Summary & Objectives
%======================================================================
\begin{abstract}
电子科技大学生物医学工程硕士在读，具备计算机视觉算法研发、端侧模型落地与跨团队工程协作经验。曾在旷视科技手机业务部门实习，独立推进注视检测模型从研发到交付，在复杂场景准确率与部署效果上取得显著提升。当前研究方向为跨设备眼动跟踪与无障碍人机交互，聚焦跨域泛化、实时交互与系统级优化，具备从算法设计到产品化实现的完整实践能力。
\end{abstract}

%======================================================================
\sectionTitle{专业技能}{\faWrench}
%======================================================================
\begin{itemize}
  \item 熟悉 Python、PyTorch、MegEngine 等深度学习开发框架，具备目标检测与目标跟踪项目经验。
  \item 熟悉 YOLOv8、ByteTrack、BoT-SORT 等视觉算法，能够完成检测/跟踪任务建模、训练与评估。
  \item 具备端侧模型优化与部署经验，熟悉 QAT/PTQ 量化流程及跨硬件平台一致性问题分析。
  \item 具备数据工程与模型工程实践能力，能够搭建数据与模型管理流程并支持多项目并行迭代。
  \item 熟悉实时眼动跟踪与人机交互系统开发，具备跨设备校准、误差优化与交互设计能力。
\end{itemize}

%======================================================================
\sectionTitle{教育背景}{\faGraduationCap}
%======================================================================
\begin{educations}
  \education%
    {2024.09}%
    [2027.06 (预计)]%
    {电子科技大学}%
    {生命科学与技术学院}%
    {生物医学工程}%
    {硕士}
  \separator{0.5ex}
  \education%
    {2020.09}%
    [2024.07]%
    {桂林电子科技大学}%
    {生命与环境科学学院}%
    {生物医学工程}%
    {学士}
\end{educations}

%======================================================================
\sectionTitle{科研经历}{\faAtom}
%======================================================================
\begin{itemize}
  \item \textbf{跨设备重定向增强方法研究} \hfill 2024.09 --- 至今
    \begin{itemize}
    \item \textbf{问题背景}: 跨设备视线估计存在尺寸差异引起的分布偏移问题，导致注视点在不同视野设备间映射不稳定。
    \item \textbf{研究内容}: 1) 提出基于重定向数据增强的跨设备泛化方法，将小视野注视数据扩增到大视野设备；2) 结合个性化注视点校准与多点快速采样策略，提升系统跨用户与跨设备鲁棒性；3) 构建跨设备评测流程，对比不同重定向策略在泛化精度与稳定性上的表现。
    \item \textbf{相关成果}: 目前已形成相关论文一篇，处于 TIP 投稿阶段。
    \end{itemize}
\end{itemize}

%======================================================================
\sectionTitle{实习经历}{\faBriefcase}
%======================================================================
\begin{experiences}
  \experience%
    [2024.04]%
    {2024.10}%
    {旷视科技 \quad 手机业务部（极感科技） \quad 算法实习生}%
    [\begin{itemize}
      \item 独立完成蔚来手机注视检测模型的研发与交付，负责算法研发与客户数据对接；基于 MegEngine 框架完成约 4 个月迭代优化，针对客户反馈 badcase 从数据与算法两侧改进。
      \item 通用场景准确率由 84\% 提升至 95\%，复杂场景准确率由 62\% 提升至 85\%，模型成功上线至蔚来最新款手机。
      \item 主导注视检测基座模型更新，使用全向人脸检测替代旧版四向检测，结合人脸 landmark 与仿射变换提升侧脸/歪脸场景鲁棒性；整体准确率提升约 2\%，半脸场景提升约 30\%。
      \item 完成算法基础设施建设与多项目支持，基于 nori2 搭建统一云端数据与模型管理流程，实现数据、模型、测试文件分离，训练效率提升约 150\%；并参与 vivo、荣耀等项目算法迭代。
      \item 参与跨算法、开发、测试、量化、研究院多部门协作，推进 QAT 与 PTQ 方案调试，解决量化与推理结果对齐问题，提升模型跨硬件平台一致性与可移植性。
    \end{itemize}]
  \separator{0.5ex}
  \experience%
    [2023.04]%
    {2023.10}%
    {欧洲蚊虫控制协会（EMCA）/ Entostudio \quad 国际实习生}%
    [\begin{itemize}
      \item 开发基于 YOLOv8 的目标检测模型，用于识别监测视频中的蟑螂与蚂蚁个体。
      \item 开发基于 ByteTrack、BoT-SORT 的目标跟踪模型，用于捕捉视频中所有个体的运动轨迹。
      \item 设计并实现基于 Qt 的软件界面，集成检测与跟踪流程，支持标定区域内目标计数与计时。
    \end{itemize}]
\end{experiences}

%======================================================================
\sectionTitle{项目经历}{\faCode}
%======================================================================
\begin{itemize}
    \item \textbf{实时眼动跟踪交互平台开发} \hfill 2024.09 --- 至今
    \begin{itemize}
      \item \textbf{项目背景}: 为提升肢体受限用户的计算机使用能力，开发基于普通摄像头的实时眼动跟踪交互平台。
      \item \textbf{核心痛点}: 纯视觉方案在实时性、注视点精度和个体差异适配方面存在明显挑战，影响可用性与稳定性。
      \item \textbf{技术方案}: 1) 构建基于笔记本摄像头的眼动跟踪与交互闭环，实现眼动驱动的鼠标控制、键盘输入及常用操作；2) 设计个性化校准与多点快速采样流程，降低跨用户误差；3) 设计多种辅助工具与交互界面，优化操作路径与反馈机制。
      \item \textbf{实验验证}: 在标准测试场景下进行时延与精度评测，系统帧率达到 30FPS 以上，平均注视点预测误差小于 2cm。
      \item \textbf{最终效果}: 平台具备稳定的人机交互能力，可支持基础输入与常用系统操作，显著提升无障碍使用体验。
      \item \textbf{技术栈}: Python, OpenCV, Qt
    \end{itemize}
\end{itemize}

%======================================================================
\sectionTitle{个人荣誉}{\faTrophy}
%======================================================================
\begin{itemize}
  \item 以第一作者发表计算机视觉领域 SCI 论文 1 篇
  \item 申请发明专利 2 项、软件著作权 1 项
  \item 第十届全国大学生生物医学工程创新设计竞赛 三等奖
  \item 第十八届“挑战杯”全国大学生课外学术科技作品竞赛 三等奖
  \item 第九届中国国际“互联网+”大学生创新创业大赛广西赛区 金奖
  \item 第三届广西大学生人工智能设计大赛 一等奖
\end{itemize}

\end{document}
